% GUID v6, Primer P10 (label box:hadron): hadron spectroscopy. Background only; it states no result of the paper.
% Place it before "Hadrons as trees of dualities" (marker %PRIMER:hadron in part1.tex): % GUID v6, Primer P10 (label box:hadron): hadron spectroscopy. Background only; it states no result of the paper.
% Place it before "Hadrons as trees of dualities" (marker %PRIMER:hadron in part1.tex): % GUID v6, Primer P10 (label box:hadron): hadron spectroscopy. Background only; it states no result of the paper.
% Place it before "Hadrons as trees of dualities" (marker %PRIMER:hadron in part1.tex): % GUID v6, Primer P10 (label box:hadron): hadron spectroscopy. Background only; it states no result of the paper.
% Place it before "Hadrons as trees of dualities" (marker %PRIMER:hadron in part1.tex): \input{primers/p10-hadron}
\begin{primer}{Hadron spectroscopy}{box:hadron}
\emph{Colour.} Quarks carry the colour charge of $SU(3)$: a quark is a triplet $\mathbf3$, an antiquark an antitriplet $\mathbf{\bar3}$.
Only colour \emph{singlets} exist as free particles, the hadrons~\cite{navas2024}. The products
$\mathbf3\otimes\mathbf{\bar3}=\mathbf1\oplus\mathbf8$ and $\mathbf3\otimes\mathbf3=\mathbf{\bar3}\oplus\mathbf6$ give the ordinary
hadrons: a meson $q\bar q$ in the singlet, and a baryon $qqq$, whose first two quarks couple to $\mathbf{\bar3}$ and then with the third to
a singlet. One-gluon exchange attracts in the singlet $q\bar q$ and, half as strongly, in the $\mathbf{\bar3}$ $qq$, and repels in the
$\mathbf8$ and the $\mathbf6$.

\emph{Diquarks.} A \emph{diquark} is a pair of quarks in the attractive $\mathbf{\bar3}$. The ``good'' diquark, with spin 0 and a flavour-
antisymmetric pair such as $[ud]$, binds most; the ``bad'' one has spin 1~\cite{jaffe2005}.

\emph{Compact or molecular.} A hadron with more than three quarks, an \emph{exotic}, can be \emph{compact}, for example a diquark and an
antidiquark bound by colour forces, $[qq][\bar q\bar q]$~\cite{maiani2005}. Or it can be a \emph{hadronic molecule}: two colour-singlet
hadrons held together by the residual strong force, like the deuteron~\cite{guo2018}. Molecules sit just below the threshold
$m_1+m_2$ of their two constituents.

\emph{The colour-spin interaction.} In the constituent-quark model a hadron's mass is the sum of effective quark masses plus a hyperfine
term proportional to $(\boldsymbol\lambda_i\cdot\boldsymbol\lambda_j)(\boldsymbol\sigma_i\cdot\boldsymbol\sigma_j)/m_im_j$ for each pair,
where the $\boldsymbol\lambda$ are colour and the $\boldsymbol\sigma$ spin matrices. This chromomagnetic interaction of De R\'ujula, Georgi
and Glashow explains, for example, why the $\Delta$ is heavier than the nucleon~\cite{derujula1975}. Karliner and Rosner (KR) refined it
with separate effective quark masses in mesons and in baryons, the baryon masses for every constituent of a compact multiquark, and binding
terms for heavy pairs with $B(QQ)=B(Q\bar Q)/2$; with it they predicted the doubly charmed baryon and doubly heavy
tetraquarks~\cite{karliner2014,karliner2017}.

\emph{Thresholds and cusps.} Where a new two-hadron channel opens, at $M=m_1+m_2$, the scattering amplitude has a branch point. A
\emph{cusp} can then appear in a mass spectrum without any state, and a state that couples to the channel is pulled towards the threshold.
Many exotic candidates lie within a few MeV of a threshold~\cite{guo2018}.

\emph{Names.} The Particle Data Group names hadrons by their quantum numbers, not by their structure~\cite{navas2024}; the
$X(3872)$, for example, is listed as $\chi_{c1}(3872)$. LHCb's convention for exotics~\cite{gershon2022} uses $T$ for four-quark and $P$
for five-quark states, with subscripts for heavy flavour, $\psi$ standing for hidden charm $c\bar c$. So $T_{cc}^+$ is $cc\bar u\bar d$,
seen just below the $D^{*+}D^0$ threshold~\cite{lhcb2022h}, and the hidden-charm pentaquark $P_c(4312)^+$~\cite{lhcb2019h} becomes
$P_\psi^N(4312)^+$.

\emph{String junctions.} In a baryon the three colour flux tubes meet at a point, a \emph{string junction}~\cite{rossiveneziano1977}. A
compact tetraquark has two junctions and a compact pentaquark three, and some models add a constant energy per
junction~\cite{karliner2017a}.
\end{primer}

\begin{primer}{Hadron spectroscopy}{box:hadron}
\emph{Colour.} Quarks carry the colour charge of $SU(3)$: a quark is a triplet $\mathbf3$, an antiquark an antitriplet $\mathbf{\bar3}$.
Only colour \emph{singlets} exist as free particles, the hadrons~\cite{navas2024}. The products
$\mathbf3\otimes\mathbf{\bar3}=\mathbf1\oplus\mathbf8$ and $\mathbf3\otimes\mathbf3=\mathbf{\bar3}\oplus\mathbf6$ give the ordinary
hadrons: a meson $q\bar q$ in the singlet, and a baryon $qqq$, whose first two quarks couple to $\mathbf{\bar3}$ and then with the third to
a singlet. One-gluon exchange attracts in the singlet $q\bar q$ and, half as strongly, in the $\mathbf{\bar3}$ $qq$, and repels in the
$\mathbf8$ and the $\mathbf6$.

\emph{Diquarks.} A \emph{diquark} is a pair of quarks in the attractive $\mathbf{\bar3}$. The ``good'' diquark, with spin 0 and a flavour-
antisymmetric pair such as $[ud]$, binds most; the ``bad'' one has spin 1~\cite{jaffe2005}.

\emph{Compact or molecular.} A hadron with more than three quarks, an \emph{exotic}, can be \emph{compact}, for example a diquark and an
antidiquark bound by colour forces, $[qq][\bar q\bar q]$~\cite{maiani2005}. Or it can be a \emph{hadronic molecule}: two colour-singlet
hadrons held together by the residual strong force, like the deuteron~\cite{guo2018}. Molecules sit just below the threshold
$m_1+m_2$ of their two constituents.

\emph{The colour-spin interaction.} In the constituent-quark model a hadron's mass is the sum of effective quark masses plus a hyperfine
term proportional to $(\boldsymbol\lambda_i\cdot\boldsymbol\lambda_j)(\boldsymbol\sigma_i\cdot\boldsymbol\sigma_j)/m_im_j$ for each pair,
where the $\boldsymbol\lambda$ are colour and the $\boldsymbol\sigma$ spin matrices. This chromomagnetic interaction of De R\'ujula, Georgi
and Glashow explains, for example, why the $\Delta$ is heavier than the nucleon~\cite{derujula1975}. Karliner and Rosner (KR) refined it
with separate effective quark masses in mesons and in baryons, the baryon masses for every constituent of a compact multiquark, and binding
terms for heavy pairs with $B(QQ)=B(Q\bar Q)/2$; with it they predicted the doubly charmed baryon and doubly heavy
tetraquarks~\cite{karliner2014,karliner2017}.

\emph{Thresholds and cusps.} Where a new two-hadron channel opens, at $M=m_1+m_2$, the scattering amplitude has a branch point. A
\emph{cusp} can then appear in a mass spectrum without any state, and a state that couples to the channel is pulled towards the threshold.
Many exotic candidates lie within a few MeV of a threshold~\cite{guo2018}.

\emph{Names.} The Particle Data Group names hadrons by their quantum numbers, not by their structure~\cite{navas2024}; the
$X(3872)$, for example, is listed as $\chi_{c1}(3872)$. LHCb's convention for exotics~\cite{gershon2022} uses $T$ for four-quark and $P$
for five-quark states, with subscripts for heavy flavour, $\psi$ standing for hidden charm $c\bar c$. So $T_{cc}^+$ is $cc\bar u\bar d$,
seen just below the $D^{*+}D^0$ threshold~\cite{lhcb2022h}, and the hidden-charm pentaquark $P_c(4312)^+$~\cite{lhcb2019h} becomes
$P_\psi^N(4312)^+$.

\emph{String junctions.} In a baryon the three colour flux tubes meet at a point, a \emph{string junction}~\cite{rossiveneziano1977}. A
compact tetraquark has two junctions and a compact pentaquark three, and some models add a constant energy per
junction~\cite{karliner2017a}.
\end{primer}

\begin{primer}{Hadron spectroscopy}{box:hadron}
\emph{Colour.} Quarks carry the colour charge of $SU(3)$: a quark is a triplet $\mathbf3$, an antiquark an antitriplet $\mathbf{\bar3}$.
Only colour \emph{singlets} exist as free particles, the hadrons~\cite{navas2024}. The products
$\mathbf3\otimes\mathbf{\bar3}=\mathbf1\oplus\mathbf8$ and $\mathbf3\otimes\mathbf3=\mathbf{\bar3}\oplus\mathbf6$ give the ordinary
hadrons: a meson $q\bar q$ in the singlet, and a baryon $qqq$, whose first two quarks couple to $\mathbf{\bar3}$ and then with the third to
a singlet. One-gluon exchange attracts in the singlet $q\bar q$ and, half as strongly, in the $\mathbf{\bar3}$ $qq$, and repels in the
$\mathbf8$ and the $\mathbf6$.

\emph{Diquarks.} A \emph{diquark} is a pair of quarks in the attractive $\mathbf{\bar3}$. The ``good'' diquark, with spin 0 and a flavour-
antisymmetric pair such as $[ud]$, binds most; the ``bad'' one has spin 1~\cite{jaffe2005}.

\emph{Compact or molecular.} A hadron with more than three quarks, an \emph{exotic}, can be \emph{compact}, for example a diquark and an
antidiquark bound by colour forces, $[qq][\bar q\bar q]$~\cite{maiani2005}. Or it can be a \emph{hadronic molecule}: two colour-singlet
hadrons held together by the residual strong force, like the deuteron~\cite{guo2018}. Molecules sit just below the threshold
$m_1+m_2$ of their two constituents.

\emph{The colour-spin interaction.} In the constituent-quark model a hadron's mass is the sum of effective quark masses plus a hyperfine
term proportional to $(\boldsymbol\lambda_i\cdot\boldsymbol\lambda_j)(\boldsymbol\sigma_i\cdot\boldsymbol\sigma_j)/m_im_j$ for each pair,
where the $\boldsymbol\lambda$ are colour and the $\boldsymbol\sigma$ spin matrices. This chromomagnetic interaction of De R\'ujula, Georgi
and Glashow explains, for example, why the $\Delta$ is heavier than the nucleon~\cite{derujula1975}. Karliner and Rosner (KR) refined it
with separate effective quark masses in mesons and in baryons, the baryon masses for every constituent of a compact multiquark, and binding
terms for heavy pairs with $B(QQ)=B(Q\bar Q)/2$; with it they predicted the doubly charmed baryon and doubly heavy
tetraquarks~\cite{karliner2014,karliner2017}.

\emph{Thresholds and cusps.} Where a new two-hadron channel opens, at $M=m_1+m_2$, the scattering amplitude has a branch point. A
\emph{cusp} can then appear in a mass spectrum without any state, and a state that couples to the channel is pulled towards the threshold.
Many exotic candidates lie within a few MeV of a threshold~\cite{guo2018}.

\emph{Names.} The Particle Data Group names hadrons by their quantum numbers, not by their structure~\cite{navas2024}; the
$X(3872)$, for example, is listed as $\chi_{c1}(3872)$. LHCb's convention for exotics~\cite{gershon2022} uses $T$ for four-quark and $P$
for five-quark states, with subscripts for heavy flavour, $\psi$ standing for hidden charm $c\bar c$. So $T_{cc}^+$ is $cc\bar u\bar d$,
seen just below the $D^{*+}D^0$ threshold~\cite{lhcb2022h}, and the hidden-charm pentaquark $P_c(4312)^+$~\cite{lhcb2019h} becomes
$P_\psi^N(4312)^+$.

\emph{String junctions.} In a baryon the three colour flux tubes meet at a point, a \emph{string junction}~\cite{rossiveneziano1977}. A
compact tetraquark has two junctions and a compact pentaquark three, and some models add a constant energy per
junction~\cite{karliner2017a}.
\end{primer}

\begin{primer}{Hadron spectroscopy}{box:hadron}
\emph{Colour.} Quarks carry the colour charge of $SU(3)$: a quark is a triplet $\mathbf3$, an antiquark an antitriplet $\mathbf{\bar3}$.
Only colour \emph{singlets} exist as free particles, the hadrons~\cite{navas2024}. The products
$\mathbf3\otimes\mathbf{\bar3}=\mathbf1\oplus\mathbf8$ and $\mathbf3\otimes\mathbf3=\mathbf{\bar3}\oplus\mathbf6$ give the ordinary
hadrons: a meson $q\bar q$ in the singlet, and a baryon $qqq$, whose first two quarks couple to $\mathbf{\bar3}$ and then with the third to
a singlet. One-gluon exchange attracts in the singlet $q\bar q$ and, half as strongly, in the $\mathbf{\bar3}$ $qq$, and repels in the
$\mathbf8$ and the $\mathbf6$.

\emph{Diquarks.} A \emph{diquark} is a pair of quarks in the attractive $\mathbf{\bar3}$. The ``good'' diquark, with spin 0 and a flavour-
antisymmetric pair such as $[ud]$, binds most; the ``bad'' one has spin 1~\cite{jaffe2005}.

\emph{Compact or molecular.} A hadron with more than three quarks, an \emph{exotic}, can be \emph{compact}, for example a diquark and an
antidiquark bound by colour forces, $[qq][\bar q\bar q]$~\cite{maiani2005}. Or it can be a \emph{hadronic molecule}: two colour-singlet
hadrons held together by the residual strong force, like the deuteron~\cite{guo2018}. Molecules sit just below the threshold
$m_1+m_2$ of their two constituents.

\emph{The colour-spin interaction.} In the constituent-quark model a hadron's mass is the sum of effective quark masses plus a hyperfine
term proportional to $(\boldsymbol\lambda_i\cdot\boldsymbol\lambda_j)(\boldsymbol\sigma_i\cdot\boldsymbol\sigma_j)/m_im_j$ for each pair,
where the $\boldsymbol\lambda$ are colour and the $\boldsymbol\sigma$ spin matrices. This chromomagnetic interaction of De R\'ujula, Georgi
and Glashow explains, for example, why the $\Delta$ is heavier than the nucleon~\cite{derujula1975}. Karliner and Rosner (KR) refined it
with separate effective quark masses in mesons and in baryons, the baryon masses for every constituent of a compact multiquark, and binding
terms for heavy pairs with $B(QQ)=B(Q\bar Q)/2$; with it they predicted the doubly charmed baryon and doubly heavy
tetraquarks~\cite{karliner2014,karliner2017}.

\emph{Thresholds and cusps.} Where a new two-hadron channel opens, at $M=m_1+m_2$, the scattering amplitude has a branch point. A
\emph{cusp} can then appear in a mass spectrum without any state, and a state that couples to the channel is pulled towards the threshold.
Many exotic candidates lie within a few MeV of a threshold~\cite{guo2018}.

\emph{Names.} The Particle Data Group names hadrons by their quantum numbers, not by their structure~\cite{navas2024}; the
$X(3872)$, for example, is listed as $\chi_{c1}(3872)$. LHCb's convention for exotics~\cite{gershon2022} uses $T$ for four-quark and $P$
for five-quark states, with subscripts for heavy flavour, $\psi$ standing for hidden charm $c\bar c$. So $T_{cc}^+$ is $cc\bar u\bar d$,
seen just below the $D^{*+}D^0$ threshold~\cite{lhcb2022h}, and the hidden-charm pentaquark $P_c(4312)^+$~\cite{lhcb2019h} becomes
$P_\psi^N(4312)^+$.

\emph{String junctions.} In a baryon the three colour flux tubes meet at a point, a \emph{string junction}~\cite{rossiveneziano1977}. A
compact tetraquark has two junctions and a compact pentaquark three, and some models add a constant energy per
junction~\cite{karliner2017a}.
\end{primer}
